\section{Recap and Basics}
% \section{Types of Goods and Markets}



\begin{frame}{Orders of Magnitude}
    \mycomment{
    Since we are dealing with large numbers (of monetary terms), we need to have some prices and figures to put them into perspective.
    
    How much is an oil tanker worth? How much its load? What is the GDP of your home country? How much is the grain worth traded per day?
    }
    \showmybitmap{../shared/bitmaps/joyofmoney.jpg}
\end{frame}

\begin{frame}{Orders of Magnitude}
    \mytable{
       \begin{tabular}{lrrll|l}
        Sc. & Acc. & IT & engl. & other (D) & \textit{example} \\
        \toprule
         $10^3$ & TEUR & K & thousand &  & \phantom{laptop/computer}\\
        $10^4$ & & 10K & &  & \\
        $10^5$ & & 100K & & \emph{lakh} & \\
        \midrule
        $10^6$ & MEUR& 1M & million & \emph{Millionen} &  \\
        $10^7$ & & 10M & &  \emph{crore} & \\
        $10^8$ & & 100M & &  & \\
        \midrule
        $10^9$ &  & 1G & billion & \emph{Milliarden}  & \\
        $10^{10}$ & & 10G & &  & \\
        $10^{11}$ &  & 100G & &  & \\
        \midrule
        $10^{12}$ &  & 1T & trillion & \emph{Billionen} & \\
        \bottomrule
    \end{tabular}
    }
\end{frame}

\begin{frame}{Orders of Magnitude}
    \mycomment{
    To have an idea about the meaning of (large, economic) numbers, it makes sense to memorise a couple of them. For instance, the GDP of your home country (and maybe of one or two larger economies) is a good standard candle. 
    
    Every now and then, it may make sense to compare big numbers to unrelated numbers. For example, there are (more than) 31 million seconds in a year, and there are 100 billion stars in the galaxy. Just how many cars can you buy for the price of one oil tanker? How many nursing homes can you build for the price of one moon mission?
    }
    \mytext{
    How much...
    \mylonglist{
    \begin{itemize}
        \item costs a car, a house?
        \item the whole super market?
        \item prime estate in London?
        \item the World Trade Center?
        \item a Panamax oil tanker, and its content?
        \item turnover in Bremerhaven per day?
        \item A container shipping Shanghai-Hamburg?
        \item The Artemis II moon mission?
        \item The companies Siemens, Saudi Aramco?
        \item US government interest payment?
        \item US government defense budget?
    \end{itemize}
    }
    }
\end{frame}

\begin{frame}{Stock Quantities vs. Flow Quantities}
    \mycomment{
    Growth rates, or rates of change in general, are often given as percentage (per unit of time, typically per year).


    }
    \mytext{
    \begin{description}
        \item[Stock quantities] describe how much there is of something at a present time. 
        \item[Flow quantities] describe how much is added or subtracted per time frame\\ (usually a year).
        \item[Orders of Magnitude] are often more important than precise numbers.
        \item[Growth Rates] are connecting stock and flow quantities; discern 1st and 2nd order derivatives to describe trends.
    \end{description}
    }
\end{frame}

\begin{frame}{Flow vs. Derivatives / Marginal Rates}
    \mycomment{
        To make numbers speak---for us, to understand numbers properly---we must always put them into perspective. So for us, growth rates are typically more interesting than absolutes; and absolutes are interesting when put into relation to other absolutes. 

        It is worth keeping in mind that `percent' is a division by one hundred. The orders of magnitude we are using, however, step up in thousands. Thus, $1\%$ of 1 Million are 10 Thousand, etc.
    }
    \showmybitmap{../shared/bitmaps/only-a-sith-deals-in-absolutes.png}
\end{frame}

\begin{frame}{Flow vs. Derivatives / Marginal Rates}
    \mycomment{Remember elasticity? Keep in mind that elasticity is a different concept to the concept of marginal rates.

    When we are looking at maxima, and that is where most equilibria in economics exist, we are seeking marginal rates of zero. So the perfect point often is where any marginal change does not change much the overall picture. 

    For example, at my optimum, I am pretty indifferent on what to spend the penny I found in my pocket.  If I am, however, very hungry, I may have a much larger marginal utility spending money on food rather than on anything else.
    }
    \mytext{
     \begin{itemize}
         \item Flow rates: how much more \textit{per period}
         \item Marginal rates: how much more \textit{per unit}
     \end{itemize}
     Mathematically, \textit{marginal rates} are the derivative of the function:
     \begin{equation}
         f'(x) = \frac{df}{dx}
     \end{equation}
     Since marginal rates are important for optimisation, they are more relevant for decision making than stock or flow values.
    }
\end{frame}

\begin{frame}{Elasticity of Demand}
    \showmytikz[1]{../shared/tikz/tikz-elasticitytwo.tex}
    \mycomment{
        So in general, a flat curve is inelastic, and a steep curve elastic. In reality, demand curves are not linear.
    
        Note that a linear demand curve does not have constant elasticity. That is because the relation the absolute values of $P,Q$ is different at different parts of the slope.

        \begin{equation}
            \eta = \frac{
                \frac{\delta Q}{Q}
            }{
                \frac{\delta P}{P}
            }
        \end{equation}

        Between $A,B$ a very small percentage change of the price $P$ has resulted in a larger percentage change of the quantity $Q$.

        At high elasticity, a tiny price change will have a large quantity effect. 

        And even a large quantity change will change the price just a little. 
        
    }
\end{frame}

\begin{frame}{Elasticity of Demand}
    \showmytikz[1,2]{../shared/tikz/tikz-elasticitytwo.tex}
    \mycomment{
         Between $C,D$ a larger percentage change of the price $P$ has resulted in a small percentage change of the quantity $Q$.

        At low elasticity, small quantity changes affect large price changes, while large price changes affect quantity little.
        
    }
\end{frame}

\begin{frame}{Elasticity of Demand}
    \showmytikz[1,2,3]{../shared/tikz/tikz-elasticitytwo.tex}
    \mycomment{
         At $U$ we have unity elasticity: a 1\% change in price will result in a 1\% change in quantity.

         It follows that maximum revenue is at the point of unity elasticity (ceteris paribus).

         Elasticity depends on the slope of the demand curve, but also on the location in the demand curve!

         Note that demand curves are in reality not linear. They may be more or less linear at the point of equilibrium, so we are typically concerned with the elasticity \textit{at that point.}
    }
\end{frame}

\begin{frame}{Elasticity}
    \showmytikz[1,2,3]{../shared/tikz/tikz-elasticity-oil.tex}
    \mycomment{
         The Iran war 2026 resulted in closure of the straight of Hormus, through which approximately $20\%$ of crude oil is transported. The oil price almost doubled. 

         Energy is the most quoted example of a commodity with very low elasticity of demand. How much less do you guess do people drive / use oil because of a $50\%$ increase in price? Well, apparently only about $20\%$ less. So here, $\eta=\frac{-0.2}{0.5}=-0.4$.

        At the time of writing (April 12th 2026) the longer term effects have not yet been visible. It is insofar not sound  to set $\Delta Q=20\%$.

        Also note that short-term elasticity and long-term elasticity differ, as buyers may be able to find substitutes after a while. 
    }
\end{frame}


\begin{frame}{Discounting}
    \mycomment{Not only do we have preferences between objects today, but we also take into account today versus tomorrow.

    For example, I can eat all my cookies today, or I can save some for tomorrow.

    While we typically do take tomorrow into account, we also typically value tomorrow somehow (a little, or much?) less than today.

    Just how much we discout tomorrow relative to today is a very good question. You can argue that the future has a (market) price, and that this price is the interest rate.

    
    }
    \mytext{
    The growth rate of anything is often more relevant than the sum. 
    \begin{equation}
        K_t = K_0 \left(1+\frac{p}{100} \right)^t
    \end{equation}

    Investors will put their money where they receive highest returns, and shy away from lower or negative returns. 
  

    }
    
\end{frame}

\begin{frame}{Rule of 72}
    \mycomment{
    Growth rates often seem small. But the the economy of a country. We have a gross domestic product in the order of magnitude of $10^{12}$ EUR (trillions). The economy grows each year by a single digit growth rate.
    \begin{itemize}
        \item By how much will the GDP grow (in the first year)? 
        \item[$\rightarrow$] $1\% $ of 1 Trillion is 10 Billion.
        \item How long will it take for the GDP to have doubled?
    \end{itemize}
    You should be able not only to compute this, but to compute an approximation in your head!
    }
    \mytext{

    If you have a growth rate of say $4\%$, how much long will it take for the stock to doube?
    
    We can use this rule of thumb for an approximation:
    \begin{equation}
        p \approx \frac{72}{\Delta t} \quad \Leftrightarrow \quad \Delta t=\frac{72}{p}
    \end{equation}

    Stock at a growth rate of 4\% will double in approximately in $\Delta t\left(p=4\right) \approx \frac{72}{4}=18$ years. 
    }
    
\end{frame}

\begin{frame}{Market Delimitation}
    \mycomment{

    Economists look at markets, that is the exchange of merchandise. Merchandise means those products are all pretty much the same, there is little differentiation by brand name etc.
    
    When we speak about a particular market, we must of course deliminate it, that is: clarifiy what we are actually talking about or what exactly we are measuring.
    }
    \mytext{
        \begin{itemize}
            \item spacial delimitation (`where')
            \item temporal delimitation (`when')
            \item by characteristics
            \begin{itemize}
                \item substitutability: buyer side
                \item substitutability: seller side
            \end{itemize}
        \end{itemize}    
        }
\end{frame}

\begin{frame}{Raw Materials, Commodities, Merchandise}
    \mycomment{
        Examples for raw materials are iron ore, and crude oil; once refined, they become commodities such as steel, and petrol; from which we produce manufactured goods, which when traded are called merchandise.

        Note that services can also be traded, and for other traded objects other vocabulary may apply. For example, you can trade bonds, or you can trade derivatives.

        Note that raw materials are a sub-class of commodities, and commodities are a sub-class of  merchandise. This is, however, not a "legal" distinction.
        }

    \showmytikz[1]{../shared/tikz/tikz-raw-commodities}
\end{frame}

\begin{frame}{Raw Materials, Commodities, Merchandise}
    \mycomment{
        Along the value chain, various means of production are needed. In principle these are capital and labour. These are in parts themselves tradable. 

    }
    \showmytikz[2]{../shared/tikz/tikz-raw-commodities}
\end{frame}

\begin{frame}{Raw Materials, Commodities, Merchandise}
    \mycomment{
        Typically, the different stages of production are not happening at the same place. So there is trade involved between most stages.
        
        Trade is not instantaneous, and it is not just magically happening. It is intrinsically linked to logistics. It includes transport over a distance, and often also over trade barriers such as customs borders or through choke points. 
    
        It also requires financial flows in the opposite direction.
    }
    \showmytikz[2,3]{../shared/tikz/tikz-raw-commodities}
\end{frame}

\begin{frame}{Value Chain, Upstream, Downstream}
    \showmytikz{../shared/tikz/tikz-valuechain.tex}
    \mycomment{
    Raw material is mined (or grown), then refined. Refined commodities are then used to produce intermediate goods, which are combined to manufactured wares. These are then branded and distributed.

    At all stages, costs occur, and value is added. Production requires not only labour, but capital in the form of machinery etc. Distribution is very labour intensive, and includes sales and marketing.    
    
    In most cases, production and distribution contribute the largest part to the value of the final consumer good. 
    }
\end{frame}

\begin{frame}{Market Characteristics}
    \mycomment{
    Every market has fundamental characteristics we must keep in mind. Above all, the size and importance of this market - for whom?
    
    Where is the merchandise traded? How is it moved? What volumes and prices are realised? Are these stable, or do they fluctuate widely? 
    
    Who are the buyers, and what do they buy for? Who are the sellers, and how do they produce the product? 
    
    Are there monopolies? 
    
    Who else is in this market: is it regulated, what role does the state / do states play? Can anybody enter this market? Are there some dominating companies? Do you need know-how, or any special privileges?

    As for the price developments and the elasticities, recent  shocks - such as the Iran war 2026 - provide a good background. How has this shock impacted the market? How have prices risen? And why?
    }
    \mytext{
        \begin{itemize}
            \item physical properties
            \item qualities / characteristics
            \item market segments
            \item volumes
            \item history
            \item adjacent markets:
            \item[] substitutes, complements
            \item players: buyers, producers, traders
            \item barriers to entry
            \item regulators, regulations
            \item factor markets
        \end{itemize}
    }
\end{frame}


\begin{frame}{Circular Model of the Economy}
    \showmytikz[1]{../shared/tikz/tikz-flowmodel-own}
    \mycomment{
        We will now look at the simplest form of markets, that is: the exchange of goods and services.
    }
\end{frame}


\begin{frame}{Public and Private Goods}
    \mycomment{
        
        Owner(s) can (by definition!) exclude non-owners from the use of the good, or allow the use under conditions such as payments. 
        
        Then the proprietor is a different person from the posessor / holder. (In English, the word "owner" may refer to either; cf. German \textit{Eigentümer} vs. \textit{Besitzer}).

        Public goods belong to everyone, like the air or common knowledge (to some degree). This is non-rival, non-excludable. But how about the water in the rivers, and the forests surrounding a village?
    }
    \mytext{
    \textbf{Private goods} are owned by individuals, but also by groups of people and legal entities. They are exclusive.
    
    \textbf{Public goods}
        belong to everyone (or nobody), like the air or common knowledge. 

    Public goods tend to be over-used (`tragedy of the commons'). 
    
    Private goods can be traded and are more likely used and allocated efficiently.
    
    }
\end{frame}

\begin{frame}{Different Types of Goods}
\mytable{
    \begin{tabular}{lll|l}
        \toprule
                    & price $\nearrow$     & income $\nearrow$    & example \\
        \midrule
                    
        \textbf{normal}      & \textbf{demand $\searrow$}  & \textbf{demand $\nearrow$}    & soap  \\
        inferior    & demand $\searrow$  & demand $\searrow$  & cheap noodles \\
        Veblen      & demand $\nearrow$    & demand $\nearrow$   & luxury watches \\
        Giffen      & demand $\nearrow$    & demand $\searrow$  & rice \\
        \bottomrule
    \end{tabular}
}
        
\end{frame}


\begin{frame}{Substitutes and Complements}
    \mycomment{Some goods stand in a relation to each other so that demand for one corresponds to demand for the other.

    Goods are complementary if they `work together', which can be simultaneously or in time. That is, if one good is part of another in the value creation chain, they will also be called complements. For example, the price of copper will correlate to the price of merchandise manufactures from copper. Commodities are economically complements of the raw material used to manufacture them!

    They are substitutes when they basically serve the same purpose, so that if one of them becomes more expensive buyers will move to the other goods. Think of different modes of transport, for instance.
    }
    \mytext{
    \begin{itemize}
        \item \textbf{substitute goods:} $A$ can be used to substitute $B$ and vice versa
        \item[] e.g. margarine and butter
        \item \textbf{complementary goods:} if $A$ is more in demand, demand for $B$ good goes also up
        \item[] e.g. printers and cartridges
        \item \textbf{independent goods:} demand for both is independent
    \end{itemize}
    Just how much one quantity changes relative to the change of another quantity is a question of elasticity.
    }
\end{frame}

